\scanpage{144}
de $E\otimes\Omega^{\unsure{\cdot|\bullet}}_{S/T}$ un complexe.

Supposons qu'il en soit ainsi. \deleted{\unsure{Rel. La $1^{\text{re}}$ tour d'un}}
\[
H^*_{\mathrm{DR}}(S,\underset{\substack{|\\\text{sous-entendu :}\\ E\text{ muni de }\delta_E\\
\text{i.e. de }d_E^0}}{E})=\mathbb H^*(S,E\otimes\Omega^{\cdot}_{S/T}). \tag*{(11)}
\]
On a une suite spectrale
\[
H^*_{\mathrm{DR}}(S,E)\Longleftarrow E_1^{pq}=
\deleted{\unsure{H^q(\mathrm{im}}}H^q(S,E\otimes\Omega^p_{S/T})) \tag*{(12)}
\]
associée à la filtration triviale du complexe $E\otimes\Omega^{\cdot}_{S/T}$.

On en \unsure{conclut|déduit} des \deleted{\unsure{aug}} homomorphismes
\[
H^*(\Gamma(S,E\otimes\Omega^{\cdot}_{S/T}))\Longrightarrow H^*_{\mathrm{DR}}(S,E) \tag*{(13)}
\]
qui sont des isom si p.ex. $S$ affine, et
\[
H^i_{\mathrm{DR}}(S,E)\longrightarrow H^i(S,E), \tag*{(14)}
\]
qui n'ont guère de raison d'être des isom. Ils sont
associés à l'hom de complexes évident (\unsure{troncature|troncation}) :
\[
E\otimes\Omega^{\cdot}_{S/T}\longrightarrow E. \tag*{(15)}
\]

\uline{NB} Pour $E$ module à conn. intégrable variable,
$H^*_{\mathrm{DR}}(S,E)$ est un $\unsure{\partial|\delta}$-foncteur en $E$ (On peut aussi
prendre la variante cat. dérivées). Les suites spectrales
et hom écrits plus haut sont fonctoriels en $E$.

\uline{Re NB} Images inverses \inserted{d'un Module à conn.} par

\diagram{\begin{tikzcd}[ampersand replacement=\&,row sep=small,column sep=small]
S\arrow[d]\& S'\arrow[l]\arrow[d]\\
T\& T'\arrow[l]
\end{tikzcd}}

carré commutatif :
tout suit \unsure{\ldots} !

En particulier ; \deleted{\unsure{img}} a) Image inverse par $T'\longrightarrow T$
(en prenant $S'=S\times_T T'$). b) Si $S\longrightarrow T'\longrightarrow T$, une
connexion rel. à $T$ en définit une rel. à $T'$, \deleted{et}
% archive sheet number: 143
